\section{Limitations}

MMLU-Redux alignment is structural, using \texttt{subject\_numeric\_index} with confidence 0.85,
not direct-id or hash-confirmed item matching. The MMLU-Redux result is weak/negative external
validation: broad/proxy diagnostics did not strongly recover the structurally aligned labels, and
subject-normalized evidence weakened the earlier raw label-error signal.

The MMLU IRT output is proxy-only and should not be described as full psychometric IRT or full 2PL.
The local Rasch/1PL optimizer is blocked for the active wide matrix because the implementation would
fit 14,080 free parameters; the 2PL output is a proxy-slope export only.
The narrow raw label-error signal is at most a preliminary review-queue cue and does not support
detection-success claims.

Synthetic validation can be circular when generators mirror detectors. The paper should therefore
report the mixed cross-flaw specificity and held-out generator results alongside
false-positive/error-control checks and materiality labels. Synthetic-harness power, numeric
calibration with confidence/logprob outputs, and human-reviewed real-benchmark threshold calibration
remain \texttt{[RESULT REQUIRED]}.

Real benchmark findings remain preliminary and protocol-scoped. The second-benchmark path is
prepared through a Kaggle notebook package, but no Kaggle outputs were imported in this run. Domain
packs are not paper-validated unless separately supported. GPQA local-panel findings are protocol
evidence, not broad validity evidence.

The subject-instability finding may depend on the model panel, HELM extraction, subject taxonomy,
and MMLU item mix. It should be treated as a reproducible panel finding, not a universal property of
all benchmarks or model rankings.
